\section{AHSME 1992}
        \begin{enumerate}
        	\item (\textit{AHSME 1992}) If $3(4x+5\pi)=P$ then $6(8x+10\pi)=$
$\text{(A) } 2P\quad \text{(B) } 4P\quad \text{(C) } 6P\quad \text{(D) } 8P\quad \text{(E) } 18P$

 \textbf{Answer: }B



\textbf{Solution}

We can see that $8x+10\pi$ is equal to $2(4x+5\pi),$ and we know that $2^2 = 4,$ so the answer is $\boxed{B}\, .$



	\item (\textit{AHSME 1992}) An urn is filled with coins and beads, all of which are either silver or gold. Twenty percent of the objects in the urn are beads. Forty percent of the coins in the urn are silver. What percent of objects in the urn are gold coins?


 $\text{(A) } 40\%\quad \text{(B) } 48\%\quad \text{(C) } 52\%\quad \text{(D) } 60\%\quad \text{(E) } 80\%$

 \textbf{Answer: }B



\textbf{Solution}

If twenty percent of the objects in the urn are beads, then the other $80$\% of the objects must be coins. Since $40$\% of the coins are silver, then $60$\% of the coins must be gold. Therefore, $60$\% of $80$\% of the objects in the urn are gold coins. Since $.6\times.8=.48$, 48 percent of the objects in the urn must be gold coins. Thus the answer is $\fbox{B}$.

 


	\item (\textit{AHSME 1992}) If $m>0$ and the points $(m,3)$ and $(1,m)$ lie on a line with slope $m$, then $m=$
$\text{(A) } 1\quad \text{(B) } \sqrt{2}\quad \text{(C) } \sqrt{3}\quad \text{(D) } 2\quad \text{(E) } \sqrt{5}$

 \textbf{Answer: }C



\textbf{Solution}

We know that the formula for slope is $m = \frac{y_2-y_1}{x_2-x_1}$We are give the points $(m,3)$ and $(1,m)$.Substituting into the slope formula, we get $\frac{m-3}{1-m}$ After taking the cross-products and solving, we get $\sqrt{3}$, which is answer choice $\fbox{C}$.

 



\textbf{Solution 2}

Using the formula for slope as above, we know that $m=\frac{m-3}{1-m}$. Multiplying both sides of this equation by $1-m$, we get that $m(1-m)=m-3$ which expands to $m-m^2=m-3$. This forms the quadratic equation $-m^2+3=0$. This means that $m^2=3$, so $m=\sqrt{3}$, which corresponds to answer choice $\fbox{C}$.

 


	\item (\textit{AHSME 1992}) If $a,b$ and $c$ are positive integers and $a$ and $b$ are odd, then $3^a+(b-1)^2c$ is


 $\text{(A) odd for all choices of c} \quad \text{(B) even for all choices of c} \quad\\ \text{(C) odd if c is even; even if c is odd} \quad\\ \text{(D) odd if c is odd; even if c is even} \quad\\ \text{(E) odd if c is not a multiple of 3; even if c is a multiple of 3}$

 \textbf{Answer: }A



\textbf{Solution}

Since 3 has no factors of 2, $3^a$ will be odd for all values of $a$. Since $b$ is odd as well, $b-1$ must be even, so $(b-1)^2$ must be even. This means that for all choices of $c$, $(b-1)^2c$ must be even because any integer times an even number is still even. Since an odd number plus an even number is odd, $3^a+(b-1)^2c$ must be odd for all choices of $c$, which corresponds to answer choice $\fbox{A}$.

 


	\item (\textit{AHSME 1992}) $6^6+6^6+6^6+6^6+6^6+6^6=$
$\text{(A) } 6^6 \quad \text{(B) } 6^7\quad \text{(C) } 36^6\quad \text{(D) } 6^{36}\quad \text{(E) } 36^{36}$

 \textbf{Answer: }B



\textbf{Solution}

We notice that $6^6+6^6+6^6+6^6+6^6+6^6=6(6^6)=\boxed{\textbf{(B) }6^7}$.

 


	\item (\textit{AHSME 1992}) If $x>y>0$ , then $\frac{x^y y^x}{y^y x^x}=$
$\text{(A) } (x-y)^{y/x}\quad \text{(B) } \left(\frac{x}{y}\right)^{x-y}\quad \text{(C) } 1\quad \text{(D) } \left(\frac{x}{y}\right)^{y-x}\quad \text{(E) } (x-y)^{x/y}$

 \textbf{Answer: }D



\textbf{Solution}

We see that this fraction can easily be factored as $\frac{x^y}{y^y}\times\frac{y^x}{x^x}$. Since $\frac{y^x}{x^x}=\frac{x^{-x}}{y^{-x}}$, this fraction is equivalent to $\left(\frac{x}{y}\right)^y\times\left(\frac{x}{y}\right)^{-x}=\left(\frac{x}{y}\right)^{y-x}\quad$, which corresponds to answer choice $\fbox{D}$.

 


	\item (\textit{AHSME 1992}) The ratio of $w$ to $x$ is $4:3$, of $y$ to $z$ is $3:2$ and of $z$ to $x$ is $1:6$. What is the ratio of $w$ to $y$?


 $\text{(A) } 1:3\quad \text{(B) } 16:3\quad \text{(C) } 20:3\quad \text{(D) } 27:4\quad \text{(E) } 12:1$

 \textbf{Answer: }B



\textbf{Solution}

$\fbox{B}$

 Converting from ratios to fractions

 $\frac{w}{x}=\frac{4}{3}$,$\frac{y}{z}=\frac{3}{2}$,$\frac{z}{x}=\frac{1}{6}$
$\frac{y}{z} \cdot \frac{z}{x} = \frac{3}{2} \cdot \frac{1}{6}$
$\frac{y}{x} = \frac{1}{4}$
$\frac{(\frac{w}{x})}{(\frac{y}{x})}=\frac{\frac{4}{3}}{\frac{1}{4}}$
$\frac{w}{y} = \frac{16}{3}$
$w$ to $y$ is $16:3$



	\item (\textit{AHSME 1992}) 
                    \begin{figure}[H]
                        \centering
                        \includegraphics[width=0.5\textwidth]{Problem/AHSME_I_1992_8_Problem_Figure_1.png}
                    \end{figure}
                    


 A square floor is tiled with congruent square tiles. The tiles on the two diagonals of the floor are black. The rest of the tiles are white. If there are 101 black tiles, then the total number of tiles is


 $\text{(A) } 121\quad \text{(B) } 625\quad \text{(C) } 676\quad \text{(D) } 2500\quad \text{(E) } 2601$

 \textbf{Answer: }E



\textbf{Solution}

$\fbox{E}$ Let the side length of the square be $s$. If $s$ is even, there are $s$ black tiles on each of the two diagonals, so the total number is $2s$. If $s$ is odd, we double-count the middle square which is on both diagonals, so the number is instead $2s-1$. In this case, since $101$ is odd, we clearly must have the second case, so $2s-1=101 \implies 2s=102 \implies s=51$, and indeed $s$ is odd as we expected. Thus the overall number of tiles is $51^2 = 2601$.

 


	\item (\textit{AHSME 1992}) 
                    \begin{figure}[H]
                        \centering
                        \includegraphics[width=0.5\textwidth]{Problem/AHSME_I_1992_9_Problem_Figure_1.png}
                    \end{figure}
                    


 Five equilateral triangles, each with side $2\sqrt{3}$, are arranged so they are all on the same side of a line containing one side of each vertex. Along this line, the midpoint of the base of one triangle is a vertex of the next. The area of the region of the plane that is covered by the union of the five triangular regions is


 $\text{(A) 10} \quad \text{(B) } 12\quad \text{(C) } 15\quad \text{(D) } 10\sqrt{3}\quad \text{(E) } 12\sqrt{3}$

 \textbf{Answer: }E



\textbf{Solution}

$\fbox{E}$First, we calculate the area of 1 triangle. For an equilateral triangle with side $s$, its area is $\frac{\sqrt{3}s^{2}}{4}$. If the side of the equilateral triangle is $2\sqrt{3}$, the area of one such triangle is $3\sqrt{3}$.There are 5 equilateral triangles in total, overlapping by an area of 4 smaller equilateral triangles. Each smaller triangle is one fourth as big as the big equilateral triangles. Therefore, we subtract the overlapping area, which is equivalent to the area of 1 big equilateral triangle. Hence, the total area is equal to the area of 4 equilateral triangles, which is $3\sqrt{3} \times 4 = 12\sqrt{3}$.

 


	\item (\textit{AHSME 1992}) The number of positive integers $k$ for which the equation

 \[kx-12=3k\]
has an integer solution for $x$ is


 $\text{(A) } 3\quad \text{(B) } 4\quad \text{(C) } 5\quad \text{(D) } 6\quad \text{(E) } 7$

 \textbf{Answer: }D



\textbf{Solution}

$\fbox{D}$
$kx -12 = 3k$
$-12=3k-kx$
$-12=k(3-x)$
$\frac{-12}{k}=3-x$

 Positive factors of $-12$:$1,2,3,4,6,12$

 6 factors, each of which have an integer solution for $x$ in $\frac{-12}{k}=3-x$



	\item (\textit{AHSME 1992}) 
                    \begin{figure}[H]
                        \centering
                        \includegraphics[width=0.5\textwidth]{Problem/AHSME_I_1992_11_Problem_Figure_1.png}
                    \end{figure}
                    


 The ratio of the radii of two concentric circles is $1:3$. If $\overline{AC}$ is a diameter of the larger circle, $\overline{BC}$ is a chord of the larger circle that is tangent to the smaller circle, and $AB=12$, then the radius of the larger circle is


 $\text{(A) } 13\quad \text{(B) } 18\quad \text{(C) } 21\quad \text{(D) } 24\quad \text{(E) } 26$

 \textbf{Answer: }B



\textbf{Solution (Similarity)}

We are given that $BC$ is tangent to the smaller circle. Using that, we know where the circle intersects $BC$, it creates a right triangle. We can also point out that since $AC$ is the diameter of the bigger circle and triangle $ABC$ is inscribed the semi-circle, that angle $B$ is a right angle. Therefore, we have $2$ similar triangles. Let's label the center of the smaller circle (which is also the center of the larger circle) as $D$. Let's also label the point where the smaller circle intersects $BC$ as $E$. So $ABC$ is similar to $DEC$. Since $DE$ is the radius of the smaller circle, call the length $x$ and since $DC$ is the radius of the bigger circle, call that length $3x$. The diameter, $AC$ is $6x$. So,

 $\frac{AB}{AC} = \frac {DE}{DC} \Rightarrow \frac{12}{6x} = \frac{x}{3x} \Rightarrow 36x = 6x^2 \Rightarrow 6 = x$

 But they are asking for the larger circle radius, so 

 $3x= \fbox{18}$
$\fbox{B}$



	\item (\textit{AHSME 1992}) Let $y=mx+b$ be the image when the line $x-3y+11=0$ is reflected across the $x$-axis. The value of $m+b$ is


 $\text{(A) -6} \quad \text{(B) } -5\quad \text{(C) } -4\quad \text{(D) } -3\quad \text{(E) } -2$

 \textbf{Answer: }C



\textbf{Solution}

First we want to put this is slope-intercept form, so we get $y=\dfrac{1}{3}x+\dfrac{11}{3}$. When we reflect a line across the x-axis, the line's x-intercept stays the same, while the y-intercept and the slope are additive inverses of the original line. Since $m+b$ is the sum of the slope and the y-intercept, we get $-\dfrac{1}{3}-\dfrac{11}{3}=\boxed{-4}$.

 


	\item (\textit{AHSME 1992}) How many pairs of positive integers $(a,b)$ with $a+b\le 100$ satisfy the equation


 
 \[\frac{a+b^{-1}}{a^{-1}+b}=13?\]
$\text{(A) } 1\quad \text{(B) } 5\quad \text{(C) } 7\quad \text{(D) } 9\quad \text{(E) } 13$

 \textbf{Answer: }C



\textbf{Solution 1}

$\fbox{C}$ We can rewrite the left-hand side as $\frac{a^2b+a}{b+ab^2}$ by multiplying the numerator and denominator by $ab$. Now we can multiply the top part of the fraction by the denominator to give $a^2b+a = 13b + 13ab^2 \implies ab(a-13b)=13b-a = -(a-13b) \implies (ab+1)(a-13b) = 0$. Now, as $a$ and $b$ are positive, we cannot have $ab+1=0$, so we must have $a-13b=0 \implies a=13b$, so the solutions are $(13,1), (26,2), (39,3), (52,4), (65,5), (78,6), (91,7)$, which is 7 ordered pairs.

 



\textbf{Solution 2}

The equation is equivalent to $\frac{a+\frac{1}{b}}{\frac{1}{a}+b}=13$ or $\frac{\frac{ab+1}{b}}{\frac{ab+1}{a}}=13$. The $ab+1$s cancel out leaving us with $\frac{a}{b}=13$. This is asking us for the positive multiples of $13$ up to $100$, and we know that $13 \cdot 7 = 91$, so our answer is $\fbox{C}$.

 


	\item (\textit{AHSME 1992}) Which of the following equations have the same graph?


 $I.\quad y=x-2 \qquad II.\quad y=\frac{x^2-4}{x+2}\qquad III.\quad (x+2)y=x^2-4$
$\text{(A) I and II only} \quad \text{(B) I and III only} \quad \text{(C) II and III only} \quad \text{(D) I,II,and III} \quad \\ \text{(E) None. All of the equations have different graphs}$

 \textbf{Answer: }E



\textbf{Solution}

The equations differ at $x=-2$. The graph of $I$ would contain the point $(-2, -4)$; $-2$ is undefined in the graph of $II$ because it gives a denominator of $0$; and the graph of $III$ contains the whole the vertical line $x = -2$ as for any $y$ value, the equation still satisfies $0 = 0$. So the answer is ${E}$.

 


	\item (\textit{AHSME 1992}) Let $i=\sqrt{-1}$. Define a sequence of complex numbers by


 
 \[z_1=0,\quad z_{n+1}=z_{n}^2+i \text{ for } n\ge1.\]
In the complex plane, how far from the origin is $z_{111}$?


 $\text{(A) } 1\quad \text{(B) } \sqrt{2}\quad \text{(C) } \sqrt{3}\quad \text{(D) } \sqrt{110}\quad \text{(E) } \sqrt{2^{55}}$

 \textbf{Answer: }B



\textbf{Solution}

$\fbox{B}$ Write out some terms: $0, i, -1+i, -i, -1+i, -i, -1+i, -i$, etc., and it keeps alternating between $-1+i$ and $-i$, so as $111$ is odd, $z_{111}$ is $-1+i$. Thus its distance from the origin is $\sqrt{(-1)^2+1^2} = \sqrt{2}$.

 


	\item (\textit{AHSME 1992}) If

 \[\frac{y}{x-z}=\frac{x+y}{z}=\frac{x}{y}\]
for three positive numbers $x,y$ and $z$, all different, then $\frac{x}{y}=$
$\text{(A) } \frac{1}{2}\quad \text{(B) } \frac{3}{5}\quad \text{(C) } \frac{2}{3}\quad \text{(D) } \frac{5}{3}\quad \text{(E) } 2$

 \textbf{Answer: }E



\textbf{Solution 1}

$\fbox{E}$ We have $\frac{x+y}{z} = \frac{x}{y} \implies xy+y^2=xz$ and $\frac{y}{x-z} = \frac{x}{y} \implies y^2=x^2-xz \implies x^2-y^2=xz$. Equating the two expressions for $xz$ gives $xy+y^2=x^2-y^2 \implies x^2-xy-2y^2=0 \implies (x+y)(x-2y)=0$, so as $x+y$ cannot be $0$ for positive $x$ and $y$, we must have $x-2y=0 \implies x=2y \implies \frac{x}{y}=2$.

 



\textbf{Solution 2}

We cross multiply the first and third fractions and the second and third fractions, respectively, for 
 \[x(x-z)=y^2\] 
 \[y(x+y)=xz\] Notice how the first equation can be expanded and rearranged to contain an $(x+y)$ term. 
 \[x^2-xz=y^2\] 
 \[x^2-y^2=xz\] 
 \[(x+y)(x-y)=xz\] We can divide this by the second equation to get 
 \[\frac{(x+y)(x-y)}{y(x+y)}=\frac{xz}{xz}\] 
 \[\frac{x-y}{y}=1\] 
 \[\frac{x}{y}-1=1\] 
 \[\frac{x}{y}=2 \rightarrow \boxed{E}\]




\textbf{Solution 3}

Since 
 \[\frac{a_1}{b_1} = \frac{a_2}{b_2} = ... = \frac{a_n}{b_n} = \frac{a_1 + a_2+ ... + a_n}{b_1 + b_2 + ... + b_n},\] we can say that 
 \[\frac{y}{x-z} = \frac{x+y}{z} = \frac{x}{y} = \frac {y+(x+y)+x}{(x-z)+z+y} = \boxed{2}.\]$\implies \boxed{(E)}$.

 $\textbf{Proof of the used property:}$

 Let 
 \[\frac{a_1}{b_1} = \frac{a_2}{b_2} = \frac{a_3}{b_3} = ... = \frac{a_n}{b_n} = r.\]

 Therefore, 
 \begin{align*} a_1 &= b_1r,\\ a_2 &= b_2r,\\ ...,\\ a_n &= b_nr \end{align*}
so 
 \[\frac{a_1 + a_2+ ... + a_n}{b_1 + b_2 + ... + b_n} = \frac{b_1r + b_2r + ... + b_nr}{b_1 + b_2 + ... + b_n}\] 
 \[= \frac{(b_1 + b_2 + ... + b_n)(nr)}{(b_1 + b_2 + ... + b_n)(n)} = r\] --- NamelyOrange

 


	\item (\textit{AHSME 1992}) The 2-digit integers from 19 to 92 are written consecutively to form the integer $N=192021\cdots9192$. Suppose that $3^k$ is the highest power of 3 that is a factor of $N$. What is $k$?


 $\text{(A) } 0\quad \text{(B) } 1\quad \text{(C) } 2\quad \text{(D) } 3\quad \text{(E) more than } 3$

 \textbf{Answer: }B



\textbf{Solution 1}

We can determine if our number is divisible by $3$ or $9$ by summing the digits. Looking at the one's place, we can start out with $0, 1, 2, 3, 4, 5, 6, 7, 8, 9$ and continue cycling though the numbers from $0$ through $9$. For each one of these cycles, we add $0 + 1 + ... + 9 = 45$. This is divisible by $9$, thus we can ignore the sum. However, this excludes $19$, $90$, $91$ and $92$. These remaining units digits sum up to $9 + 1 + 2 = 12$, which means our units sum is $3 \pmod 9$. As for the tens digits, for $2, 3, 4, \cdots , 8$ we have 

 


	\item (\textit{AHSME 1992}) The increasing sequence of positive integers $a_1,a_2,a_3,\cdots$ has the property that


 
 \[a_{n+2}=a_n+a_{n+1} \text{  for all } n\ge 1.\]


 If $a_7=120$, then $a_8$ is


 $\text{(A) } 128\quad \text{(B) } 168\quad \text{(C) } 193\quad \text{(D) } 194\quad \text{(E) } 210$

 \textbf{Answer: }D



\textbf{Solution}

$\fbox{D}$ Let $a_{1} = a, a_{2} = b$, so $5a + 8b = 120$. Now $8b$ and $120$ are divisible by $8$, so $5a$ is divisible by 8, so $a$ is divisible by 8. It's now easy to try the multiples of $8$ to get that $a = 8, b=10$ (all the other possibilities violate the condition $a < b$, which comes from the fact that the sequence is increasing). Hence $a_8 = 8a + 13b = 8 \times 8 + 13 \times 10 = 194$.

 


	\item (\textit{AHSME 1992}) For each vertex of a solid cube, consider the tetrahedron determined by the vertex and the midpoints of the three edges that meet at that vertex. The portion of the cube that remains when these eight tetrahedra are cut away is called a cubeoctahedron. The ratio of the volume of the cubeoctahedron to the volume of the original cube is closest to which of these?


 $\text{(A) } 75\%\quad \text{(B) } 78\%\quad \text{(C) } 81\%\quad \text{(D) } 84\%\quad \text{(E) } 87\%$

 \textbf{Answer: }D



\textbf{Solution}

$\fbox{D}$ Let the cube have side length 1, and place the cube in the coordinate plane. Then we can pick any vertex, get the coordinates of the midpoints, and hence find the three vectors that define the tetrahedron (the vectors from the chosen vertex to each midpoint). Now using $\frac{1}{6}\|a.(b\times c)\|$, we can find the volume of one of the tetrahedra, then multiply it by 8, and subtract from 1 to get $\frac{5}{6}$, which is closest to $84\%$.

 


	\item (\textit{AHSME 1992}) 
                    \begin{figure}[H]
                        \centering
                        \includegraphics[width=0.5\textwidth]{Problem/AHSME_I_1992_20_Problem_Figure_1.png}
                    \end{figure}
                    
Part of an "n-pointed regular star" is shown. It is a simple closed polygon in which all $2n$ edges are congruent, angles $A_1,A_2,\cdots,A_n$ are congruent, and angles $B_1,B_2,\cdots,B_n$ are congruent. If the acute angle at $A_1$ is $10^\circ$ less than the acute angle at $B_1$, then $n=$
$\text{(A) } 12\quad \text{(B) } 18\quad \text{(C) } 24\quad \text{(D) } 36\quad \text{(E) } 60$

 \textbf{Answer: }D



\textbf{Solution}

If we sum up the angles to obtain 360, we can see that the B angles add to the sum and the A angles subtract from the sum (an easy way of looking at this is by using the opposing angle theorem: if A[n] = B[n] than their total contribution is 0). Thus we have B[1] + B[2] + ... + B[n] - A[1] - A[2] - ... A[n] = 360. But every pair of A,B has a total 'angle contribution' of 10, thus there are 36 $\fbox{D}$  pairs of A,B.

 


	\item (\textit{AHSME 1992}) For a finite sequence $A=(a_1,a_2,...,a_n)$ of numbers, the \textit{Cesáro sum} of A is defined to be 
$\frac{S_1+\cdots+S_n}{n}$ , where $S_k=a_1+\cdots+a_k$ and $1\leq k\leq n$. If the Cesáro sum of
the 99-term sequence $(a_1,...,a_{99})$ is 1000,  what is the Cesáro sum of the 100-term sequence 
$(1,a_1,...,a_{99})$?


 $\text{(A) } 991\quad \text{(B) } 999\quad \text{(C) } 1000\quad \text{(D) } 1001\quad \text{(E) } 1009$

 \textbf{Answer: }A



\textbf{Solution 1}

Let us define the Cesáro total of a particular sequence to be n * Cesáro sum. We can see that the Cesáro total is 
 \[S_1 + S_2 + S_3 + ... + S_n\]
 \[= a_1 + (a_1 + a_2) + (a_1 + a_2 + a_3) + \ldots + (a_1 + a_2 + \ldots + a_n)\]
 \[= n \cdot a_1 + (n - 1) \cdot a_2 + \ldots + a_n.\]If we take this to be the Cesáro total for the second sequence, we can see that all the terms but the first term make up the first sequence. Since we know that the Cesáro total of the original sequence is $1000 * 99 = 99,000,$ than the Cesáro total of the second sequence is $n \cdot a_1 + 99,000 = 100 \cdot 1 + 99,000 = 99,100.$ Thus the Cesáro sum of the second sequence is $\frac{99,100}{100} = \boxed{991, A}\, .$




\textbf{Solution 2}

We know that the Cesáro sum of the first $99$ numbers is 

 
 \[\frac{a_1+a_1+a_2+a_1+a_2+a_3+...+a_{99}}{99}=100.\]

 Multiplying by $99$, we see that 

 
 \[a_1+a_1+a_2+a_1+a_2+a_3+...+a_{99}=100\cdot{99}.\]

 Now, we append $1$ to the list of $a_n,$ so our new sum becomes 

 
 \[\frac{1+1+a_1+1+a_1+a_2+...+a_{99}}{100} \rightarrow\]
 \[\frac{100+a_1+a_1+a_2+...+a_{99}}{100}.\]

 We've estabished that $a_1+a_1+a_2+a_1+a_2+a_3+...+a_{99}=100\cdot{99}$, so

 
 \[\frac{100+a_1+a_1+a_2+...+a_{99}}{100}=\frac{100+99\cdot{100}}{100}.\]

 Dividing everything by $100$, the sum of the sought sequence is $1+990=991 \rightarrow \boxed{A}$



	\item (\textit{AHSME 1992}) Ten points are selected on the positive $x$-axis,$X^+$, and five points are selected on the positive $y$-axis,$Y^+$. The fifty segments connecting the ten points on $X^+$ to the five points on $Y^+$ are drawn. What is the maximum possible number of points of intersection of these fifty segments that could lie in the interior of the first quadrant?


 $\text{(A) } 250\quad \text{(B) } 450\quad \text{(C) } 500\quad \text{(D) } 1250\quad \text{(E) } 2500$

 \textbf{Answer: }B



\textbf{Solution}

$\fbox{B}$ We can pick any two points on the $x$-axis and any two points on the $y$-axis to form a quadrilateral, and the intersection of its diagonals will thus definitely be inside the first quadrant. Hence the answer is $\binom{10}{2} \times \binom{5}{2} = 450$.

 


	\item (\textit{AHSME 1992}) Let $S$ be a subset of $\{1,2,3,...,50\}$ such that no pair of distinct elements in $S$ has a sum divisible by $7$. What is the maximum number of elements in $S$?


 $\text{(A) } 6\quad \text{(B) } 7\quad \text{(C) } 14\quad \text{(D) } 22\quad \text{(E) } 23$

 \textbf{Answer: }E



\textbf{Solution}

The fact that $x \equiv 0 \mod 7 \Rightarrow 7 \mid x$ is assumed as common knowledge in this answer.

 First, note that there are $8$ possible numbers that are equivalent to $1 \mod 7$, and there are $7$ possible numbers equivalent to each of $2$-$6 \mod 7$.

 Second, note that there can be no pairs of numbers $a$ and $b$ such that $a \equiv -b$ mod $7$, because then $a+b | 7$. These pairs are $(0,0)$, $(1,6)$, $(2,5)$, and $(3,4)$. Because $(0,0)$ is a pair, there can always be $1$ number equivalent to $0 \mod 7$, and no more.

 To maximize the amount of numbers in S, we will use $1$ number equivalent to $0 \mod 7$, $8$ numbers equivalent to $1$, and $14$ numbers equivalent to $2$-$5$. This is obvious if you think for a moment. Therefore the answer is $1+8+14=23$ numbers.$\fbox{E}$



	\item (\textit{AHSME 1992}) Let $ABCD$ be a parallelogram of area $10$ with $AB=3$ and $BC=5$. Locate $E,F$ and $G$ on segments $\overline{AB},\overline{BC}$ and $\overline{AD}$, respectively, with $AE=BF=AG=2$. Let the line through $G$ parallel to $\overline{EF}$ intersect $\overline{CD}$ at $H$. The area of quadrilateral $EFHG$ is


 $\text{(A) } 4\quad \text{(B) } 4.5\quad \text{(C) } 5\quad \text{(D) } 5.5\quad \text{(E) } 6$

 \textbf{Answer: }C



\textbf{Solution 1}

$\fbox{C}$ Use vectors. Place an origin at $A$, with $B = p, D = q, C = p + q$. We know that $\|p \times q\|=10$, and also $E=\frac{2}{3}p, F=p+\frac{2}{5}q, G = \frac{2}{5}q$, and now we can find the area of $EFHG$ by dividing it into two triangles and using cross-products (the expressions simplify using the fact that the cross-product distributes over addition, it is anticommutative, and a vector crossed with itself gives zero).

 



\textbf{Solution 2}

We note that $ABFG$ is a parallelogram because $AG = BF = 2$ and $AG \parallel BF$. Using the same reasoning, $GFCD$ is also a parallelogram. 

 
Assume that the height of parallelogram $ABFG$ with respect to base $AB$ is $x$. Then, the area of parallelogram $ABFG$ is $AB * x$. The area of triangle $EFG$ is $\frac{AB * x}{2}$, which is half of the area of parallelogram $ABFG$.

 
Likewise, the area of triangle $FGH$ is half the area of parallelogram $GFCD$.

 
Thus, $[EFHG] = [EFG] + [FGH] = 1/2[ABFG] + 1/2[GFCD] = 1/2[ABCD] = 1/2(10) = \boxed{5}$




\textbf{Solution 3}

After drawing and labeling your diagram, connect line $GF$ and notice that there are now two triangles $EGF$ and $GFH$. $ABFG$ and $GFCD$ are both parallelograms with a triangle inscribed in them with one of the sides as bases, so the combined area of the triangles, which forms $EFHG$, is just the area over $2$, so the answer is just $\frac{10}{2}=\boxed{5}$



	\item (\textit{AHSME 1992}) In $\triangle{ABC}$, $\angle{ABC}=120^\circ,AB=3$ and $BC=4$. If perpendiculars constructed to $\overline{AB}$ at $A$ and to $\overline{BC}$ at $C$ meet at $D$, then $CD=$
$\text{(A) } 3\quad \text{(B) } \frac{8}{\sqrt{3}}\quad \text{(C) } 5\quad \text{(D) } \frac{11}{2}\quad \text{(E) } \frac{10}{\sqrt{3}}$

 \textbf{Answer: }E



\textbf{Solution 1 (Extending Line Segments)}

We begin by drawing a diagram.
                    \begin{figure}[H]
                        \centering
                        \includegraphics[width=0.5\textwidth]{Solution/AHSME_I_1992_25_Solution_Figure_1.png}
                    \end{figure}
                    We extend $CB$ and $DA$ to meet at $E.$ This gives us a couple right triangles in $CED$ and $BEA.$
                    \begin{figure}[H]
                        \centering
                        \includegraphics[width=0.5\textwidth]{Solution/AHSME_I_1992_25_Solution_Figure_2.png}
                    \end{figure}
                    We see that $\angle E = 30^\circ$. Hence, $\triangle BEA$ and $\triangle DEC$ are 30-60-90 triangles.

 Using the side ratios of 30-60-90 triangles, we have $BE=2BA=6$. This tells us that $CE=BC+BE=4+6=10$. Also, $EA=3\sqrt{3}$.

 Because $\triangle DEC\sim\triangle BEA$, we have 
 \[\frac{10}{3\sqrt{3}}=\frac{CD}{3}.\]Solving the equation, we have
 \begin{align*} \frac{CD}3&=\frac{10}{3\sqrt{3}}\\ CD&=3\cdot\frac{10}{3\sqrt{3}}\\ CD&=\frac{10}{\sqrt{3}}\ \end{align*}
Hence, $CD=\boxed{\textbf{E}}$.

 



\textbf{Solution 2 (Cyclic Quadrilaterals, Right Triangles, LoC)}

Since $\angle{A}+\angle{C} = 180^{\circ}, ABCD$ is cyclic. Using Ptolemy's Theorem gets $\text{(1): } 4AD + 3CD = \sqrt{37}BD.$

 Right triangles $\Delta ABD$ and $\Delta CBD$ obtain $\text{(2): } 9+AD^2=BD^2$ and $\text{(3):} 16+CD^2=BD^2,$ respectively.

 Seeing squares in $\text{(2)}$ and $\text{(3)}$, we square $\text{(1)}$  and get $\text{(4): } 16AD^2+9CD^2 + 24AD\cdot CD = 37BD^2.$

 We don't like that $AD\cdot CD$ term, but fortunately LoC exists: $37 = AD^2 + CD^2 - 2*AD*CD*\cos(60^{\circ})$. Solving for $AD\cdot CD$ and plugging it into $\text{(4)}$, and using $AD^2 = 7 + CD^2$ and $BD^2 = 16 + CD^2$ from the first two equations, gets $40(CD^2+7)+33CD^2 - 24\cdot 37 = 37(16+CD^2).$

 Solve for $CD = \boxed{\textbf{E}}$.

 


	\item (\textit{AHSME 1992}) 
                    \begin{figure}[H]
                        \centering
                        \includegraphics[width=0.5\textwidth]{Problem/AHSME_I_1992_26_Problem_Figure_1.png}
                    \end{figure}
                    


 Semicircle $\widehat{AB}$ has center $C$ and radius $1$. Point $D$ is on $\widehat{AB}$ and $\overline{CD}\perp\overline{AB}$. Extend $\overline{BD}$ and $\overline{AD}$ to $E$ and $F$, respectively, so that circular arcs $\widehat{AE}$ and $\widehat{BF}$ have $B$ and $A$ as their respective centers. Circular arc $\widehat{EF}$ has center $D$. The area of the shaded "smile" $AEFBDA$, is


 $\text{(A) } \left(2-\sqrt{2}\right)\pi\quad \text{(B) } 2\pi-\pi \sqrt{2}-1\quad \text{(C) } \left(1-\frac{\sqrt{2}}{2}\right)\pi\quad\\ \text{(D) } \frac{5\pi}{2}-\pi\sqrt{2}-1\quad \text{(E) } \left(3-2\sqrt{2}\right)\pi$

 \textbf{Answer: }B



\textbf{Solution}

$\fbox{B}$ The area of the entire outer shape is the area of sector $ABE$, plus the area of sector $ABF$, minus the area of triangle $ABD$ (since it is part of both sectors), plus the area of sector $DEF$. We know $AC = CD = 1$, so the sector angles for $ABE$ and $ABF$ are $45$ degrees, and the radius of both of them is $2$. The radius of $DEF$ is $DE = BE - BD = 2 - BD$, and $BD$ can be found using Pythagoras in triangle $BCD$, giving $BD = \sqrt{2}$ and $DE = 2 - \sqrt{2}$, so after doing all the calculations, the area of the entire outer shape is $\pi(\frac{5}{2} - \sqrt{2}) - 1$. To get the area of the smile, we need to subtract the area of semicircle $ABD$, which is $\frac{1}{2} \pi 1^2 = \frac{\pi}{2}$, so the answer is $\pi(\frac{5}{2} - \frac{1}{2} - \sqrt{2}) - 1$ = $2\pi - \pi \sqrt{2} - 1$.

 


	\item (\textit{AHSME 1992}) A circle of radius $r$ has chords $\overline{AB}$ of length $10$ and $\overline{CD}$ of length 7. When $\overline{AB}$ and $\overline{CD}$ are extended through $B$ and $C$, respectively, they intersect at $P$, which is outside of the circle. If $\angle{APD}=60^\circ$ and $BP=8$, then $r^2=$

 \textbf{Answer: }D



\textbf{Solution}


                    \begin{figure}[H]
                        \centering
                        \includegraphics[width=0.5\textwidth]{Solution/AHSME_I_1992_27_Solution_Figure_1.png}
                    \end{figure}
                    

 Applying Power of a Point on $P$, we find that $PC=9$ and thus $PD=16$. Observing that $PD=2BP$ and that $\angle BPD=60^{\circ}$, we conclude that $BPD$ is a $30-60-90$ right triangle with right angle at $B$. Thus, $BD=8\sqrt{3}$ and triangle $ABD$ is also right. Using that fact that the circumcircle of a right triangle has its diameter equal to the hypotenuse, we compute using the Pythagorean Theorem $AD=2r=2\sqrt{73}$. From here we see that $r^2=73$. The answer is thus $\fbox{D}$.

 


	\item (\textit{AHSME 1992}) Let $i=\sqrt{-1}$. The product of the real parts of the roots of $z^2-z=5-5i$ is


 $\text{(A) } -25\quad \text{(B) } -6\quad \text{(C) } -5\quad \text{(D) } \frac{1}{4}\quad \text{(E) } 25$

 \textbf{Answer: }B



\textbf{Solution 1}

Applying the quadratic formula gives 
 \[z=\frac{1\pm\sqrt{21-20i}}{2}\]

 Let 
 \[\sqrt{21-20i}=a+bi\] where $a$ and $b$ are real. Squaring both sides and equating real and imaginary parts gives 
 \[a^2-b^2=21\] 
 \[2ab=-20\] Substituting $b=-\frac{10}{a}$, letting $a^2=n$ and solving gives $n=25, -4$, from which we see that $a=5$ and $b=-2$.

 Replacing the square root in the quadratic formula with $5-2i$ and simplifying gives the two roots $3-i$ and $-2+i$. The product of their real parts is $\boxed{\textbf{(B)}\ {-6}}$.

 



\textbf{Solution 2}

Complete the square to get $\left(z-\frac{1}{2}\right)^2 = \frac{21-20i}{4}$.  

 Since $\left|\frac{21-20i}{4}\right| = \frac{29}{4}$, and $\frac{21-20i}{4} = \frac{29}{4}\left(\frac{21}{29}-\frac{20}{29}i\right)$, $\left|\frac{21}{29}-\frac{20}{29}i\right| = 1$.  Therefore, there exists $\theta$ such that $0 \leq \theta \leq 2\pi$ such that $\cos \theta = \frac{21}{29}$ and $\sin \theta = -\frac{20}{29}$.

 So $\left(z-\frac{1}{2}\right)^2 = \frac{29}{4}(\cos \theta + i \sin \theta)$, and thus $z = \frac{1}{2} \pm \frac{\sqrt{29}}{2}\left(\cos\frac{\theta}{2} + i \sin\frac{\theta}{2}\right)$.

 Since $\cos \theta > 0$ and $\sin \theta < 0$, $\frac{3\pi}{2} \leq \theta \leq 2\pi$.  Thus $\frac{3\pi}{4} \leq \frac{\theta}{2} \leq \pi$, so $\cos \frac{\theta}{2} < 0$ and $\sin \frac{\theta}{2} > 0$. 

 Therefore, $\cos \frac{\theta}{2} = -\sqrt{\frac{1+\cos\theta}{2}} = -\frac{5}{\sqrt{29}}$ and $\sin \frac{\theta}{2} = \sqrt{\frac{1-\cos\theta}{2}} = \frac{2}{\sqrt{29}}$.

 So $z = \frac{1}{2} \pm \frac{\sqrt{29}}{2}\left(-\frac{5}{\sqrt{29}} + \frac{2}{\sqrt{29}}i\right) = \frac{1}{2} \pm \left(-\frac{5}{2} + i\right)$.  Therefore, the two solutions are $z = 3 - i$ and $z = -2 + i$, and the product of their real parts is $\boxed{\textbf{(B)}\ -6}$.

 


	\item (\textit{AHSME 1992}) An "unfair" coin has a $2/3$ probability of turning up heads. If this coin is tossed $50$ times, what is the probability that the total number of heads is even?


 $\text{(A) } 25\left(\frac{2}{3}\right)^{50}\quad \text{(B) } \frac{1}{2}\left(1-\frac{1}{3^{50}}\right)\quad \text{(C) } \frac{1}{2}\quad \text{(D) } \frac{1}{2}\left(1+\frac{1}{3^{50}}\right)\quad \text{(E) } \frac{2}{3}$

 \textbf{Answer: }D



\textbf{Solution}

Doing casework on the number of heads (0 heads, 2 heads, 4 heads...), we get the equation 
 \[P=\left(\frac{1}{3} \right)^{50}+\binom{50}{2}\left(\frac{2}{3} \right)^{2}\left(\frac{1}{3} \right)^{48}+\dots+\left(\frac{2}{3} \right)^{50}\] This is essentially the expansion of $\left(\frac{2}{3}+\frac{1}{3} \right)^{50}$ but without the odd power terms. To get rid of the odd power terms in $\left(\frac{2}{3}+\frac{1}{3} \right)^{50}$, we add $\left(\frac{2}{3}-\frac{1}{3} \right)^{50}$ and then divide by $2$ because the even power terms that were not canceled were expressed twice. Thus, we have 
 \[P=\frac{1}{2}\cdot\left(\left(\frac{1}{3}+\frac{2}{3} \right)^{50}+\left(\frac{2}{3}-\frac{1}{3} \right)^{50} \right)\] Or 
 \[\frac{1}{2}\left(1+\left(\frac{1}{3} \right)^{50} \right)\] which is equivalent to answer choice $\fbox{D}$.

 



\textbf{Solution 2}

Denote $P_n$ as the probability that the number of heads is even. Then we have two cases:

 $\bold{Case 1}$: (There are an odd number of heads for $(n-1)$ coin flips).Well to solve this case, we first ask the question; "What is the probability we get an odd number of heads for $(n-1)$ flips?" That is precisely $1-P_{n-1}$. Now, since we want an even number of heads, what is the probability that on the $n-th$ coin flip we get a head? $\frac{2}{3}$. So we have $\frac{2}{3}(1-P_{n-1})$.

 $\bold{Case 2}$: (There are an even number of heads for $(n-1)$ coin flips). Since we now already have an even number of heads, we now want to flip and get a tail, or else the number of heads for $n$ flips will be odd. :/ So what is the probability we get a tail? $1-\frac{2}{3}=\frac{1}{3}$. Thus we have $\frac{1}{3}(P_{n-1})$. 

 Adding both cases, we have that $P_n=\frac{2}{3}-\frac{1}{3}P_{n-1}$. Now doing some tedious casework because wolfram alpha cannot find a pattern, we get the same answer as in $\boxed {D)}$. 

 Sorry for the very hand-wavy and non-rigorous answer! At least mine has a cool recurrence relation! :D

 (Or just come up with the general formula $P_n=\frac{1}{2}(1+\frac{1}{3^n})$ Simple observation really. If you really want to get mathy, prove this through induction).

 



\textbf{Solution 3}

Obviously there is a higher chance of getting an even amount of heads, because from $0$ to $50$ there are $26$ even numbers and $25$ odd numbers. Keep this in mind while we cross out answer choices.Answer choice A is way too small.B is less than $\frac{1}{2}$, which can't be it.C is exactly $\frac{1}{2}$, so it can't be it either.D works.E is way too large.

 Thus the answer is $\boxed{D}$



	\item (\textit{AHSME 1992}) Let $ABCD$ be an isosceles trapezoid with bases $AB=92$ and $CD=19$. Suppose $AD=BC=x$ and a circle with center on $\overline{AB}$ is tangent to segments $\overline{AD}$ and $\overline{BC}$. If $m$ is the smallest possible value of $x$, then $m^2$=


 $\text{(A) } 1369\quad \text{(B) } 1679\quad \text{(C) } 1748\quad \text{(D) } 2109\quad \text{(E) } 8825$

 \textbf{Answer: }B



\textbf{Solution}

Note that the center of the circle is the midpoint of $AB$, call it $M$. When we decrease $x$, the limiting condition is that the circle will eventually be tangent to segment $AD$ at $D$ and segment $BC$ at $C$. That is, $MD\perp AD$ and $MC\perp BC$. 

 From here, we drop the altitude from $D$ to $AM$; call the base $N$. Since $\triangle DNM \sim \triangle ADM$, we have
 \[\frac{DM}{19/2}=\frac{46}{DM}.\]Thus, $DM=\sqrt{19\cdot 23}$. Furthermore, $x^2=AM^2-DM^2=46^2-19\cdot 23=1679, \boxed{B}.$



\end{enumerate}